\documentclass[11pt,a4paper]{article}
\usepackage[T1]{fontenc}
\usepackage{lmodern,microtype}
\usepackage[margin=2.4cm]{geometry}
\usepackage{amsmath,amssymb,amsthm,mathrsfs}
\usepackage{booktabs,tabularx,longtable,array,calc}
\usepackage{enumitem,xcolor,needspace}
\usepackage{titling}
\setlength{\droptitle}{-1.5\baselineskip}
\PassOptionsToPackage{hyphens}{url}
\usepackage[hidelinks]{hyperref}
\emergencystretch=2em
\providecommand{\tightlist}{\setlength{\itemsep}{0pt}\setlength{\parskip}{0pt}}
\theoremstyle{plain}
\newtheorem{maintheorem}{Theorem}
\renewcommand{\themaintheorem}{\Alph{maintheorem}}
\theoremstyle{definition}
\newtheorem{cited}{Cited result}
\newenvironment{numbered}[3]{\par\addvspace{\medskipamount}\Needspace{3\baselineskip}\noindent
  \textbf{#1 #2}\if\relax\detokenize{#3}\relax\else\ (#3)\fi\textbf{.}\ \itshape\ignorespaces}
  {\par\addvspace{\smallskipamount}}
\newenvironment{numberedup}[3]{\par\addvspace{\medskipamount}\Needspace{3\baselineskip}\noindent
  \textbf{#1 #2}\if\relax\detokenize{#3}\relax\else\ (#3)\fi\textbf{.}\ \ignorespaces}
  {\par\addvspace{\smallskipamount}}
\newcommand{\uses}[1]{\par\nopagebreak[4]\noindent{\small\textsf{Uses:} #1}\par\smallskip}
\newcommand{\step}[2]{\par\smallskip\noindent\textbf{Step #1} (#2).\ \ignorespaces}
\newcommand{\by}[1]{{\small[#1]}}
\newcolumntype{L}[1]{>{\raggedright\arraybackslash}p{#1}}
\newcolumntype{Y}{>{\raggedright\arraybackslash}X}
\newcommand{\fixedhead}[1]{\section*{#1}}
\title{Polynomial-time unit-job scheduling\\for every fixed number of machines $m\geq3$}
\author{Human contributor: entered at deposit}
\date{15 November 2026 \quad Version 1 \quad CC0 1.0}
\begin{document}
\maketitle
\vspace{-1.5\baselineskip}

\begin{abstract}
For every fixed $m\geq3$, unit-job scheduling with explicitly listed precedence constraints on $m$ identical machines admits exact deadline decision and optimal-schedule construction in $O_m((L+2)^{60m^2+129m+2})$ deterministic multitape steps; a parsed, padded deadline of size $N=mT$ costs $O_m(N^{60m^2+129m+1})$.
OpenAI proved the three-machine result (2026); drbh supplied four-machine Lean code (2026); Jansen states the fixed-machine result in arXiv:2611.09392 (2026).
This generalization gives a self-contained proof with leaf budget $12m+25$ and an itemized tape-cost derivation, without a priority claim.
A finite dynamic program uses Boolean descriptions; an injection into one boundary slot bounds the lists, and local restriction identities control direction changes.
GPT-6 Astra, ChatGPT Pro mode, produced the proof and the current step-by-step, computational and source audit. Three earlier reports are by Claude Opus 5.5, effort max, Message Batches API. Human understanding: none.
Source-access and archival qualifications are recorded in Limits and open points.
\end{abstract}

\section{Result}
\subsection{Statement}

The input explicitly lists a DAG $G=(V,E)$, $n=|V|\geq1$ jobs, and binary integers $m\geq3,T$.
A schedule is a map $\tau:V\to\{1,\ldots,T\}$ with at most $m$ jobs per slot and $\tau(u)<\tau(v)$ on every arc.
The makespan is $\max\tau$. The canonical encoding and its bit length $L$ are specified in Section~9.

\begin{cited}[OpenAI, {\cite[Theorem~1.1]{P}}]\label{cit:base}
For $m=3$, deadline decision and witness construction for $1\leq T\leq n$, and optimum-schedule construction, take polynomial deterministic multitape time in the explicit input length.
\end{cited}
\uses{Source statement opened and compared in source check S-P; the proof below is self-contained.}

\noindent\begin{tabularx}{\linewidth}{@{}lYll@{}}
\toprule
ID & Result & Status & Proof\\
\midrule
A & Exact algorithm for every fixed $m\geq3$ & proved & 9.1--9.3\\
B & Hierarchy with $12m+25$ global leaves & proved & 8.5\\
\bottomrule
\end{tabularx}

\enlargethispage{\baselineskip}
\begin{maintheorem}[Theorem 1.1 of the checked return]\label{thm:A}
One deterministic algorithm, given the input above, decides deadline feasibility, constructs a feasible schedule on acceptance, and constructs an optimal schedule. For each fixed $m$ all these tasks use
\[
 O_m\bigl((L+2)^{60m^2+129m+2}\bigr)\tag{1.1}
\]
steps on a deterministic Turing machine with a fixed number of tapes.
For $1\leq T<n$ and $n\leq mT$, its deadline computation on the parsed, canonically numbered padded DAG with $N=mT$ jobs, including its witness, uses
\[
 O_m\bigl(N^{60m^2+129m+1}\bigr).\tag{1.2}
\]
Proposition 9.2 and Corollary 9.3 give explicit effective coefficient bounds $A(m,12m+25)$ and $G(m)$, respectively.
\end{maintheorem}
\uses{Theorem 9.1; Proposition 9.2; Corollary 9.3.}
\pagebreak

\begin{maintheorem}[Theorem 8.5 of the checked return]\label{thm:B}
For $m\geq3$, every full schedule of the padded job set $J$, $|J|=mT$, admits a description hierarchy as in Definition 2.9. The root has its one-leaf description $J$, every other node has a global description of at most $12m+25$ leaves, and every marked original-time prefix has a relative description of at most $6m+16$ leaves, independently of depth.
\end{maintheorem}
\uses{Definition 2.9; Proposition 7.3; Lemmas 8.2--8.4.}

\noindent\textbf{Instance.} For $m=4$, $K_4=73$, (1.2) has exponent $1477$, and (1.1) has exponent $1478$.
\textbf{Not claimed:} a bound polynomial in $m$; a practical algorithm; extensions to release dates, individual deadlines, unequal machine speeds or unequal job lengths; formalization of this proof.

\subsection{Status}
\noindent\begin{tabularx}{\linewidth}{@{}L{3.0cm}Y@{}}
\toprule
Kind & generalization\\
Claim & proved\\
Checks & Current text: F, D, C, S by GPT-6 Astra, ChatGPT Pro mode, 15 November 2026; verdict: proved as stated. Supplied return: three earlier reads by Claude Opus 5.5, effort max, Message Batches API.\\
Human understanding & none\\
Formalization & none for this proof\\
Open points & No mathematical gap found in Sections 2--9; source-access and archival qualifications are listed below.\\
\bottomrule
\end{tabularx}

\subsection{Known and new}
%%KNOWN_TABLE%%

\subsection{Proof map}

\noindent\textbf{Idea.} Pad to full $m$-job slots and enumerate bounded Boolean descriptions of job sets. A state path joins smaller accepted sets through single slots. Separators split a feasible schedule into a hierarchy. The map of Lemma 5.5 sends surviving injection entries injectively into one boundary slot; the restricted equalities of Lemma 6.1 keep their descriptions bounded through direction changes.

\begin{enumerate}[nosep]
\item Foundations, finite search and its certificate: Lemmas 2.1--2.7, Proposition 2.8, Definition 2.9 and Lemma 2.10.
\item Exchanges, separators and walks: Lemmas 3.1--3.6 and Corollary 3.7.
\item Centers, sides and child labels: Lemmas 4.1--4.5.
\item Carried lists and the capacity injection: Lemmas 5.1--5.5.
\item Direction changes and formula counts: Lemmas 6.1--6.5.
\item Exterior invariants and induction: Lemmas 7.1--7.2 and Proposition 7.3.
\item Global child identities, prefixes and the budget: Lemmas 8.1--8.4 and Theorem 8.5.
\item Correctness and itemized multitape costs: Theorem 9.1, Proposition 9.2 and Corollary 9.3; compact decision input: Proposition 9.4.
\end{enumerate}

\noindent\textbf{Hard steps.} Lemma 5.5 proves the $m$-entry injection bound; Lemmas 6.2 and 6.5 prevent leaf counts from growing with depth; Lemma 7.1 propagates exterior qualifiers through the first advance using (I3). Their independent derivations are in the current audit. The constants are
\[
 A_m=3m+(3m+7)+2=6m+9,\qquad
 K_m=(A_m+2)+(A_m+5)=12m+25.
\]

\noindent\textbf{Reading routes.} The result and model are in Section 1.1 and Section 9. Readers comparing with the three-machine proof can follow 3.3, 5.5, 6.2, 6.5, 7.1, 8.4 and 9.2. The full proof runs through Sections 2--9 in order; Proposition 7.3 checks the construction's induction order.

\noindent\textbf{External results used.} None load-bearing. Source checks compare the construction with OpenAI's preprint and the course; every needed mathematical step is proved here. The external results in the preceding table provide historical and algorithmic context.

\noindent\textbf{Dependencies.} The following table covers all 44 numbered statements of Sections 2--9, including Definition 2.9. The source column is checked against the pinned public source of [P]. Its entries identify the three-machine counterparts; the last column records the change proved here. Definitions and displayed constructions supplement the numbered dependencies where the Uses line says so.

%%DEPENDENCY_TABLE%%

%%NOTATION_TABLE%%

\setcounter{section}{1}
